\section{Implications for Reversible-Circuit Sampling}

The transposition model is a clean abstraction of reversible identity-circuit
sampling.  In a general reversible gate library, identity circuits are closed
walks in the Cayley graph of the generated permutation group.  The fixed-length
uniform target is again a bridge distribution.  Easy samplers such as
$CC^{-1}$, adjacent cancellations, or products of relators produce identity
circuits efficiently, but they are structurally biased.

Theorem~\ref{thm:dense-rank-recursion} shows that this obstacle can be overcome
in the all-transpositions model at a dense length: one can sample a bridge
negligibly close to uniform in time polynomial in the permutation degree and
word length.  This is a genuinely global recursive construction, not a mixing
claim for Hurwitz or local heat-bath dynamics.  Its use of the stabilizer chain
$S_2\subset\cdots\subset S_N$, star transpositions, and explicit symmetric-group
spectra is specialized; it does not automatically transfer to an arbitrary
reversible gate library.

The transposition results illustrate the discipline needed for less biased
identity-circuit sampling:
\begin{itemize}
\item define the target distribution explicitly;
\item prove the actual output law, whether by detailed balance or direct
  recursive disintegration;
\item test exact small cases before scaling;
\item track invariants that local moves may preserve accidentally;
\item separate exact endpoint validity, statistical approximation, and
  computational pseudorandomness.
\end{itemize}

For cryptographic use, conditioning on identity means the output permutation is
known.  Any pseudorandomness claim must therefore concern the circuit
description or the difficulty of distinguishing the generated identity word
from a chosen target distribution over identity words.  The present work does
not prove such a cryptographic claim.  It supplies a statistically natural
reference sampler in one abstract gate model; converting that into a
cryptographic construction still requires an efficiently represented gate
library and an explicit hardness assumption about circuit descriptions.
